% 取消下一行注释可生成 handout（无逐步显示）；保持注释则生成正常 slide。
% \PassOptionsToClass{handout}{beamer}
\documentclass[Main-Prob-All.tex]{subfiles}
\begin{document}

\title[概率论]{第十六讲：数学期望}
%\author[张鑫 {\rm Email: xzhangseu@seu.edu.cn} ]{\large 张 鑫}
\institute{\large \textrm{Email: xzhangseu@seu.edu.cn} \\ \quad  \\
	\vspace{0.3cm}
	%\trc{公共邮箱: \textrm{zy.prob@qq.com}\\
		%	\hspace{-1.7cm}  密 \qquad 码: \textrm{seu!prob}}
}
\date{}








% \begin{CJK*}{GBK}{song}

%   \title{简单随机抽样}
%   \author{林语堂}
%   \institute{University}
%   \date{2009-03-31}
%   \date{}
%   \begin{frame}
%     \titlepage
%   \end{frame}

%   \begin{frame}
%     \frametitle{第二章：简单随机抽样}
%     \tableofcontents
%   %     You might wish to add the option[pausesections]
%   \end{frame}

%   \AtBeginSubsection[]{
%   \begin{frame}<beamer>
%     \frametitle{Outline}
%     \tableofcontents[currentsection,currentsubsection]
%   \end{frame}
% }
%   定义目录页


{ \setbeamertemplate{footline}{}
  \begin{frame}
    \titlepage
  \end{frame}
}



\section{随机变量的数字特征与特征函数}
\subsection{数学期望}
\begin{frame}
	\frametitle{平均值与加权平均值}
	有甲乙两名射手，其射击技术可用下表表出 \begin{table}
		甲射手 \qquad \qquad\qquad \qquad \qquad \qquad 乙射手 \\
		\vspace{0.4cm}
		\rowcolors[]{1}{blue!20}{blue!10}
		\begin{tabular}{c|c|c|c}
			\hline
			\rowcolor{blue!50}
			击中环数  &$8$ & $9$&$10$\\
			\hline
			概 \quad 率 & 0.3 & 0.1  & 0.6 \\
			\hline
		\end{tabular}
		\begin{tabular}{c|c|c|c}
			\hline
			\rowcolor{blue!50}
			击中环数  &$8$ & $9$&$10$\\
			\hline
			概 \quad 率 & 0.2 & 0.5  & 0.3 \\
			\hline
		\end{tabular}
	\end{table}
	试问哪一个射手技术较好？

	\pause
	\begin{itemize}[<+-|alert@+>]
		\item 显然这个问题的答案不是一眼就可以看出来的；
		\item 这也表明分布列虽然完整的描述了随机变量，但却不够集中的反应其变化情况；
		\item 有必要找一些量来更集中，更概括的描述随机变量；
		\item 所要找的量多是某种平均值.
	\end{itemize}
\end{frame}

\begin{frame}
	\frametitle{平均值与加权平均值}
	\begin{itemize}[<+-|alert@+>]
		\item 最常见的平均值求法:$\bar{x}=\dfrac{x_1+x_2+\cdots+x_n}{n}$;
		\item 没有考虑每个数据相对重要性，比如一小学生考试成绩为：语文 95 分，数学 85 分，常识 60 分，若按上述计算其平均成绩为 $\bar{x}=80$;
		\item 上述计算没有考虑到三个科目的相对重要性，不太能反映学生的真正成绩：如果在这个年级中，每周有语文 10 节课，数学 8 节课，常识 2 节课，则利用下述的平均计算其平均成绩似科更合理一些:
		\begin{eqnarray*}
			\bar{x}_w=95*\dfrac{10}{20}+85*\dfrac{8}{20}+60*\dfrac{2}{20}=87.5
		\end{eqnarray*}
		\item 加权平均：给定权 $w_i$ 满足 $\sum_{i=1}^nw_i=1$, 则
		\begin{eqnarray*}
			\bar{x}_w=\sum_{i=1}^nw_ix_i
		\end{eqnarray*}
		称为加权平均值.
	\end{itemize}
\end{frame}

\begin{frame}
	有甲乙两名射手，其射击技术可用下表表出 \begin{table}
		甲射手 \qquad \qquad\qquad \qquad \qquad \qquad 乙射手 \\
		\vspace{0.4cm}
		\rowcolors[]{1}{blue!20}{blue!10}
		\begin{tabular}{c|c|c|c}
			\hline
			\rowcolor{blue!50}
			击中环数  &$8$ & $9$&$10$\\
			\hline
			概 \quad 率 & 0.3 & 0.1  & 0.6 \\
			\hline
		\end{tabular}
		\begin{tabular}{c|c|c|c}
			\hline
			\rowcolor{blue!50}
			击中环数  &$8$ & $9$&$10$\\
			\hline
			概 \quad 率 & 0.2 & 0.5  & 0.3 \\
			\hline
		\end{tabular}
	\end{table}
	\pause 考虑以概率为权重的加权平均:
	\begin{eqnarray*}
		\bar{x}_{\mbox{甲}}=8*0.3+9*0.1+10*0.6=9.3\\
		\pause \bar{x}_{\mbox{乙}}=8*0.2+9*0.5+10*0.3=9.1
	\end{eqnarray*}
	\pause 则平均起来，甲每枪射中 9.3 环，乙每枪射中 9.1 环，故甲的技术更好一些.
\end{frame}





\begin{frame}
	\frametitle{数学期望的定义}
	\begin{defi}
		设离散型随机变量的分布列为 \begin{eqnarray*}
			\left(\begin{array}{ccccc}
				x_1,&x_2, &\cdots, &x_k, &\cdots\\
				p_1,&p_2, &\cdots, &p_k, &\cdots
			\end{array}\right)
		\end{eqnarray*}
		如果
		\begin{eqnarray*}
			\sum_{i=1}^\infty|x_i|p_i<\infty
		\end{eqnarray*}
		则称
		\begin{eqnarray*}
			E(X)=\sum_{i=1}^\infty x_iP(X=x_i)=\sum_{i=1}^\infty x_ip_i
		\end{eqnarray*}
		为随机变量 $X$ 的数学期望，或称该分布的数学期望，简称期望或均值。若级数 $\sum_{i=1}^\infty|x_i|p_i$ 不收敛，则称 $X$ 的期望不存在.
	\end{defi}
\end{frame}
\begin{frame}
	\frametitle{从离散到连续}
	\begin{itemize}[<+-|alert@+>]
		\item 假设连续型随机变量 $X$ 的分布密度为 $p (x)$;
		\item 考虑随机变量 $X$ 的如下近似：记 $A_i:=\{\omega:X (\omega)\in (x_i, x_{i+1}]\}$
		\begin{eqnarray*}
			\tilde{X}:=\sum_{i}x_iI_{A_i}(\omega)
		\end{eqnarray*} 其中 $I_{A_i}(\omega)=\left\{
		\begin{array}{ll}
			1,& \omega\in A_i\\
			0,& \omega\notin A_i
		\end{array}\right.
		$.

		显然 $\tilde{X}$ 为离散型随机变量，且
		\[P(\tilde{X}=x_i)=\int_{x_i}^{x_{i+1}}p(x)dx\approx p(x_i)(x_{i+1}-x_i);\]
		\item 故 $\tilde{X}$ 的期望为
		\begin{eqnarray*}
			E(\tilde{X})\approx \sum_i x_ip(x_i)(x_{i+1}-x_i)\rightarrow \int xp(x)dx
		\end{eqnarray*}

	\end{itemize}
\end{frame}

\begin{frame}
	\frametitle{连续型随机变量数学期望的定义}
	\begin{defi}
		设连续随机变量 $X$ 的密度函数为 $p (x)$. 如果
		\begin{eqnarray*}
			\int_{-\infty}^{\infty}|x|p(x)dx<\infty,
		\end{eqnarray*}
		则称
		\begin{eqnarray*}
			E(X):=\int_{-\infty}^\infty xp(x)dx
		\end{eqnarray*}
		为随机变量 $X$ 的数学期望，简称期望或均值。若 $\int_{-\infty}^\infty|x|p (x) dx$ 不收敛，则称 $X$ 的期望不存在.
	\end{defi}
\end{frame}

\begin{frame}
	\frametitle{数学期望的一般情形}
	\begin{itemize}[<+-|alert@+>]
		\item 若随机变量 $X$ 的分布函数为 $F (x)$;
		\item 类似于连续情形，考虑随机变量 $X$ 的如下近似:
		\begin{eqnarray*}
			\tilde{X}:=\sum_{i}x_iI_{A_i}(\omega)
		\end{eqnarray*}
		\vspace{-0.3cm} 其中 $A_i:=\{\omega:X (\omega)\in (x_i, x_{i+1}]\}, I_{A_i}(\omega)=\left\{
		\begin{array}{ll}
			1,& \omega\in A_i\\
			0,& \omega\notin A_i
		\end{array}\right.
		$.

		显然 $\tilde{X}$ 为离散型随机变量，且
		\[P(\tilde{X}=x_i)=P(X\in (x_i,x_{i+1}])=F(x_{i+1})-F(x_i);\]
		\item 故 $\tilde{X}$ 的期望为
		\begin{eqnarray*}
			E (\tilde{X})=\sum_i x_i\left[F (x_{i+1})-F (x_i)\right]\rightarrow \int xdF (x) \quad \textcolor{red}{\mbox{斯蒂尔切斯积分}}
		\end{eqnarray*}

	\end{itemize}
\end{frame}

\begin{frame}
	\frametitle{数学期望的一般定义}
	\begin{defi}
		假设随机变量 $X$ 的分布函数为 $F (x)$, 则定义
		\begin{eqnarray*}
			E(X)=\int_{-\infty}^{+\infty}xdF(x)
		\end{eqnarray*}
		为随机变量 $X$ 的数学期望。这里我们要求上述积分绝对收敛，否则称数学期望不存在.
	\end{defi}


\end{frame}

\begin{frame}
	\frametitle{斯蒂尔切斯 (Stieltjes) 积分的性质}
	\vspace{-0.6cm}
	\begin{eqnarray*}
		I=\int_{-\infty}^{+\infty}g(x)dF(x)
	\end{eqnarray*}
	\pause  \begin{itemize}[<+-|alert@+>]
		\item 当 $F (x)$ 为右连续的阶梯函数，在 $x_i (i=1,2,\cdots,)$ 具有跃度 $p_i$ 时，上面的积分化为
		\begin{eqnarray*}
			I=\sum_{i}g(x_i)\bigg(F(x_i)-F(x_{i-1})\bigg)=\sum_i g(x_i)p_i
		\end{eqnarray*}
		\item 当 $F (x)$ 存在导数 $F'(x)=p (x)$ 时，上述积分化为普通积分
		\begin{eqnarray*}
			I=\int_{-\infty}^{+\infty}g(x)p(x)dx
		\end{eqnarray*}
		\item 线性性质
		{\small\begin{eqnarray*}
				\int_{-\infty}^{+\infty}(ag_1(x)+bg_2(x)dF(x)=a\int_{-\infty}^{+\infty}g_1(x)dF(x)+b\int_{-\infty}^{+\infty}g_2(x)dF(x)
		\end{eqnarray*}}
		\item 若 $g (x)\ge 0$, $F (x)$ 单调不减，$b>a$, 则 $\int_a^bg (x) dF (x)\ge 0$
	\end{itemize}
\end{frame}

\begin{frame}
	\frametitle{几点注记}
	\begin{itemize}
		\item   根据斯蒂尔切斯 (Stieltjes) 积分的定义可证:%在引入上述相关的积分后，
		\begin{eqnarray*}
			\int_{(a, b]}dF(y)&=&F(b)-F(a)\\
			\int_{(-\infty, x]}dF(y)&=&\lim_{a\rightarrow-\infty}\int_{(a,x]}dF(y)=\lim_{a\rightarrow-\infty}[F(x)-F(a)]\\
			&=&F(x)=P(X\le x)\\
			\int_BdF(y)&=&P(X\in B)
		\end{eqnarray*}

		\item 数学期望的另一计算公式：关于概率测度的积分
		{\small\begin{eqnarray*}
				\textcolor{blue}{E (X)}&=& \pause \textcolor{blue}{\int_{-\infty}^\infty xdF (x)} \pause \quad \textcolor{red}{\mbox{ 注意到} F (x)=P (X (\omega)\le x):=P (X^{-1}(x))}\\
				&\xlongequal[]{x=X(\omega)}&\pause \int_{\Omega}X(\omega)dP(X^{-1}(X(\omega)))=\pause \textcolor{blue}{\int_{\Omega}X(\omega)dP(\omega)}
		\end{eqnarray*}}

	\end{itemize}

\end{frame}
\end{document}
